\scanpage{180}
$X$ \deleted{plat} au dessus de $S$

$Y$ sous-schéma divisoriel de $X$ transversal aux fibres est défini par un \uline{faisceau inversible} $\underline L$ sur $X$ et une section $s$ de $\underline L$ \deleted{\unsure{transv.}} non div de $0$ sur les fibres

Partons d'un faisceau inversible $\underline L$ sur $X$, et déterminons les \deleted{\unsure{faisc|faisceaux}} sous-schémas divisoriels $Y$ de $X$ transversaux aux fibres qui définissent des $\underline L(Y)$ \deleted{\unsure{iso}} qui sont \deleted{\unsure{Pic} $Y$-isomorphes}\deleted{\unsure{équivalents}} \deleted{à $\underline L$} isomorphes [resp\textsuperscript{t.} $\unsure{S|Y}$-équivalents, resp\textsuperscript{t.} \deleted{$Y$-équivalents} Pic\deleted{\unsure{$X/S$}}-équivalents rel. à $Y$] à $\underline L$. // Les premiers sont ceux qui se définissent par des sections de $\underline L$ non diviseurs de $0$ sur les fibres, sections qu'il faut prendre mod $\Gamma(S,f_*(\underline O_X^*))$ ; si $f_*(\underline O_X)=\underline O_S$, on les prend donc mod $\underline O_S^*$, donc on trouve les \unsure{quotients|partie} \deleted{de $\underline H_f^0(\underline L$} d'une \inserted{certaine} partie $\underline H^0_f(\underline L)^*$ de $\underline H_f^0(\underline L)$ par le groupe $\underline H^0_f(\underline O_X^*)=\unsure{\underline H_f^0|\underline W_f^0}(\underline O_X)^*$, donc dans le cas favorable par $\underline O_S^*$. --- Dans le deuxième cas, on trouve les \deleted{\unsure{éléments}}\inserted{sections} de $\underline H^0_f(\underline L)^*/\underline H_f^0(\underline O_X)^*$ donc dans le cas favorable de $\underline H^0_f(\underline L)^*/\underline O_S^*$. ---
% archive sheet number: 179
